% !TeX root = ../main.tex

\chapter{绪论}

\section{研究背景与意义}

近年来，大语言模型（Large Language Model，LLM）在自然语言处理与代码任务中取得了显著进展。以Transformer~\cite{vaswani2017attention}架构为基础的大模型在文本生成、代码补全、程序理解、测试生成与缺陷修复等任务上展现出较强能力，推动了软件工程研究从“模型直接生成代码”逐步迈向“模型驱动复杂软件工程流程”的新阶段。特别是随着GPT-4~\cite{achiam2023gpt}、Gemini~\cite{team2023gemini}、Llama~\cite{touvron2023llama}以及DeepSeek-Coder~\cite{guo2024deepseekcoder}等模型的发展，大模型在代码场景中的应用边界不断扩展，越来越多研究开始将其与工具调用、环境交互和多步规划机制结合，构建能够自主完成复杂任务的代码智能体（Code Agent）。

在各类软件工程任务中，自动程序修复（Automated Program Repair，APR）是代码智能体最具代表性的应用方向之一。与传统仅面向单文件或局部片段的代码生成任务不同，真实软件仓库中的修复任务往往需要智能体理解自然语言问题描述，在大规模代码仓库中搜索相关文件与函数，结合跨文件依赖关系生成补丁，并在真实环境中运行测试验证修复效果。这种问题求解过程本质上不再是单次输入输出，而是一条由搜索、阅读、编辑、执行与验证构成的多轮交互链路。因此，真实代码修复任务对大模型提出的要求，已经从“能否生成代码”转变为“能否在复杂仓库环境中高效完成问题求解”。

这一任务演进与近年的评测基准发展高度一致。围绕代码智能体的研究，社区逐步从HumanEval、MBPP等以函数级代码生成为主的基准，扩展到CrossCodeEval、RepoQA等更强调跨文件上下文与仓库级理解的任务，并进一步发展到SWE-bench~\cite{jimenez2024swebench}、OmniGIRL~\cite{guo2025omnigirlmultilingualmultimodalbenchmark}等更接近真实GitHub issue resolution的问题场景。尤其是OmniGIRL从多语言、多模态GitHub issue resolution的角度强调了真实软件工程任务的复杂性，而SWE-bench则进一步将这一问题具体化为仓库级代码修复与测试验证任务，为研究真实代码修复智能体提供了标准化评测环境。这些工作共同表明，代码智能体的研究正在从局部代码生成走向真实仓库级问题解决，而后者天然伴随着更长的上下文、更复杂的交互链路以及更高的推理成本。

然而，随着代码智能体能力的增强，一个长期未被充分重视的问题日益突出，即\textbf{推理成本过高}。在真实仓库级修复任务中，智能体往往需要多次搜索代码、重复读取文件、频繁调用执行环境，并在每轮交互中保留历史执行轨迹与工具返回结果。随着执行轮次的增加，输入给模型的上下文不断累积，token 消耗持续增长。对于基于 API 的模型调用场景而言，高 token 消耗不仅直接转化为经济成本，还会带来更高的响应时延、更差的并发能力和更低的工程部署可行性。换言之，即便一个代码修复智能体能够成功完成任务，若其代价过高，也难以在真实软件工程场景中得到有效应用。

值得注意的是，这里的成本问题并不能简单等同为“最终 prompt 太长”。对于多轮交互式代码智能体而言，成本更准确地说是一种\textbf{过程性成本}。它来源于整个任务求解过程中每一次模型调用的输入与输出累积，受到搜索范围、代码阅读量、历史 observation 冗余、失败重试链路等多种因素共同影响。此前关于长上下文代码压缩的研究表明，代码任务中的压缩并非简单的文本缩减问题，压缩比例提升往往会伴随任务质量下降，这说明代码场景中的上下文优化具有更强的结构敏感性与任务依赖性。因此，在代码修复智能体中，若仅从表层 prompt 压缩出发，而忽视搜索、阅读与历史上下文管理等过程性因素，往往难以真正解决高成本问题。

基于上述背景，本文聚焦于代码修复智能体在真实仓库级修复任务中的高推理成本问题，旨在回答两个层面的核心挑战：第一，代码修复智能体的推理成本究竟来自哪些环节；第二，如何在不显著损害修复能力的前提下降低其推理成本。围绕这一问题，本文以 \texttt{mini-swe-agent} 为执行框架，以 \texttt{SWE-bench} 为主要实验基准，构建完整的实验与分析闭环，对代码修复智能体的执行轨迹进行阶段化分析，并在此基础上提出面向代码理解阶段的上下文压缩与成本优化思路。该研究不仅有助于提升代码智能体在真实软件工程场景中的可用性，也为大语言模型驱动的软件工程自动化研究提供了新的分析视角。


\section{研究问题提出}

尽管近年来围绕代码智能体与自动程序修复的研究持续增加，但现有工作大多关注修复成功率、补丁正确性、基准排名或工具框架设计，对于推理成本这一问题的系统研究仍然相对不足。尤其是在仓库级代码修复场景中，任务复杂度、上下文规模和多轮交互深度显著提升，传统面向单轮输入的分析方式已难以准确刻画真实开销。基于此，本文围绕以下四个核心问题展开研究。

\textbf{研究问题一：代码修复智能体的推理成本主要来自哪些环节？} 真实代码修复任务通常包括问题理解、代码定位、上下文组织、补丁生成、测试验证与多轮调整等过程，不同阶段对 token 消耗的贡献并不一致。若无法明确主要成本来源，后续优化便缺乏针对性。

\textbf{研究问题二：仓库级修复任务为何比普通代码生成更容易引发上下文膨胀？} 相较于单文件补全或函数级代码生成，仓库级修复需要处理更大的搜索空间、更复杂的跨文件依赖以及不断累积的执行历史。理解这些因素如何共同作用，是构建有效成本分析框架的前提。

\textbf{研究问题三：已有上下文压缩与检索增强方法在代码修复智能体中是否真正适用？} 现有长上下文压缩工作与检索增强方法，往往主要面向代码补全、代码问答或通用文本任务设计，未必能够直接适应 repair agent 的多轮执行链路、结构化代码信息与历史 observation 管理需求。评估其适配性，有助于明确后续设计方向。

\textbf{研究问题四：如何在保证合理修复能力的前提下降低代码修复智能体的推理成本？} 在真实部署中，成本与能力之间往往存在权衡关系。本文希望探索面向仓库级修复任务的低成本执行路径，使代码智能体能够在保持可接受修复成功率的同时，提高效率并降低开销。


\section{研究目标与研究内容}

围绕上述研究问题，本文以代码修复智能体在仓库级修复任务中的高推理成本问题为研究核心，结合前期在长上下文代码任务和 GitHub issue resolution 场景方面的研究基础，开展以下四方面工作。

\textbf{（1）构建面向代码修复智能体的实验与分析闭环。} 以 \texttt{mini-swe-agent} 为执行主体，以 \texttt{SWE-bench} 为主要任务场景，搭建包含 issue 输入、仓库环境构建、补丁生成、轨迹记录与 token 统计的完整实验流程，为后续成本分析提供可复现的数据基础。

\textbf{（2）建立代码修复智能体的阶段化成本分析框架。} 从任务执行过程而非单轮 prompt 的视角出发，对代码修复智能体的运行轨迹进行阶段划分，分析不同阶段、不同轮次以及不同样本类型的 token 消耗特征，识别主要高成本来源与潜在的低收益模式。

\textbf{（3）提出面向仓库级修复任务的上下文压缩与成本优化思路。} 结合前期关于长上下文代码压缩的经验，围绕检索范围缩减、历史 observation 压缩、无效循环控制等方向，提出更适合代码修复智能体的成本优化策略。

\textbf{（4）讨论代码修复智能体中的成本—能力权衡关系。} 在方法比较与实验分析的基础上，联合考察 token 消耗、修复成功率、运行时延与单位成功成本，为代码智能体的工程部署与后续研究提供依据。


\section{论文结构安排}

本文共分七章，各章内容安排如下。

\textbf{第一章：绪论。} 本章介绍研究背景与意义，阐明代码修复智能体高推理成本问题的研究价值，提出核心研究问题，说明研究目标与研究内容，并概述全文结构。

\textbf{第二章：相关背景。} 本章介绍本文涉及的核心背景知识，包括大语言模型与代码任务、代码智能体的基本工作机制、自动程序修复与仓库级修复任务、GitHub issue resolution 与 \texttt{SWE-bench} 场景，以及 token、上下文窗口与推理成本等基本概念。

\textbf{第三章：相关工作与研究基础。} 本章回顾基于大语言模型的自动程序修复、代码智能体、仓库级代码任务、长上下文代码压缩与成本优化等相关研究，并在此基础上说明本文与前期已发表/已投稿工作的关系，明确本文的研究基础与切入点。

\textbf{第四章：问题定义与成本分析框架。} 本章对代码修复智能体的任务流程进行形式化抽象，定义推理成本及其组成形式，分析仓库级修复中高成本问题的来源，并给出本文的成本分析框架。

\textbf{第五章：实验平台与成本优化方法。} 本章介绍实验平台的设计与实现，说明轨迹记录与成本统计方式，并给出面向代码修复智能体的上下文压缩与成本优化方法。

\textbf{第六章：实验结果与分析。} 本章介绍实验设置，对不同方法在代码修复任务中的成本与效果进行比较，分析阶段化成本结构，讨论失败样本中的高成本模式以及成本—能力权衡关系。

\textbf{第七章：总结与展望。} 本章总结全文的主要研究内容与结论，讨论本文工作的局限性，并展望未来可能的研究方向。
